The key challenge of affective computing is to translate a subjective emotional experience into an objective state.
Most recent advances in that field have relied on facial expressions as indicators for inner emotional states.

Two main approaches have emerged to accomplish this task: the dimensional and
the categorical approach. The former describes an emotion along a set of independent dimensions, such as valence, arousal,
and dominance. This is mathematically elegant, but suffers from lacking consensus on how
to represent more complex emotions. The categorical approach uses discrete emotional states
that are often defined by their corresponding facial expression. This approach has been
popularized through the widespread use of emojis that show such expressions and which humans can instantly recognize.
However, their categorical nature make it difficult to compute numerical calculations and statistical measures.

This paper seeks to bring the power of facial visualization to the dimensional approach.
It introduces a new dimensional model and a corresponding graphical representation.
According to five emotional dimensions (valence, arousal, dominance, contempt and control)
a unique visualization with strong emotional content is derived. Their comical style is free of references to gender, age or race.
Each such face could be an emoji, but they are parametric and allow all intermediate states to be represented.
To test their effectiveness I hired untrained crowd workers to rate the visualizations
along the five emotional dimensions. The results confirm that the encoded dimensions can indeed be
recorgnized intuitively and without further training.